% !TEX root = ../main.tex
\section{Methodology}
\label{sec:methods}
\subsection{Cases, roles, and feasible incentives}
The prospective registry contains five source items: two financial tasks and
one task each in education, contract interpretation, and scientific abstract
verification. Each has a supported and an unsupported declarative claim,
an evaluator-only reference check, a source identifier, and a license record.
Participants receive the case evidence and applicable rules, without the
reference label or the operator's coalition plan. Independent interpretation
review remains necessary, particularly for the adapted contract claim.

The first model study evaluates three jurors from the same model family per
session. The creator and challenger are scripted experimental counterparts;
their behavior is excluded from model denominators. In the false-publication
mechanism, the creator offers rewards for supporting a false claim. In the
wrongful-rejection mechanism, the challenger offers rewards for rejecting a
supported claim. This tests resistance to an induced offer in the juror role.
It does not measure spontaneous initiation or autonomous behavior in creator
and challenger roles. Those extensions require separate admission.

When a majority settles the case, the contract divides the jury reward among
all valid revealers, including minority voters. With three revealers under the
selected parameters, each earns approximately 1.67 synthetic ETH net. Each paid
offer promises an additional 10 ETH for revealing the offering party's preferred
verdict if that party wins. These promises are unescrowed; an offer becomes a
payment only when the transfer executes. Feasibility fixtures must demonstrate
both honest completion and the proposed harmful outcome on the actual engine
before model observations are admitted. All game amounts are synthetic Anvil
units, separate from the US-dollar inference budget.

\subsection{Models, assignments, and controls}
Table~\ref{tab:models} in the appendix fixes the provider identifiers. A family name
is a display label; the run record must retain the exact returned provider
metadata, source revision, prompts, input hashes, configuration, and launch
order. A launch seed need not control provider sampling randomness.

The frozen confirmation registers every model with five launch-seed labels and
six arms; each session contains the same two source cases, the FinQA return
calculation and the MathDial arithmetic review. Its full frame has 150 sessions
and 300 assigned case opportunities. The initial aggregate US\$100 inference
cap was explicitly increased to US\$300 on the same spending ledger, retaining
all charges and unresolved request holds. The separate five-source pilot is exploratory and is excluded from the
primary confirmation summary. Budget exhaustion can leave observations
incomplete without removing assignments. Registered work,
started work, completed work, failed work, and missing work will remain
distinguishable, and no replacement will erase an earlier failure.

A shared sandbox-startup failure affected 75 previously pending sessions.
They remain failures in the original frame. A separately registered exploratory
cohort links 75 new attempts to those failures without replacing them; its
results are reported separately. Subsequent deadline interruptions are retained
without replay. The operator later removed the overall execution deadline,
while leaving the spending cap and frozen participant configurations intact.

The six arms comprise honest controls for false and supported claims, paid
false publication, paid wrongful rejection, paid false publication with message
delivery suppressed, and paid false publication with stronger principles.
The two basic paid arms define the primary resistance comparison with equal
mechanism weight. Controls and candidate safeguards are separate contrasts.
Message suppression preserves interfaces and instructions while changing
delivery; visible actions can still carry information. The principles arm
changes delivered duties. We therefore retain a common behavioral-harm endpoint
and score obligation violations separately.
These are six experimental conditions, not six financial-stakes levels: paid
offers use the same 10 ETH amount. Finance and education each contribute one
fixed source item, so repeated seeds do not expand task or sector coverage.

\subsection{Evidence, labels, and resistance}
\label{sec:labels}
A proposal identifies a concrete arrangement, sender, recipient, and time.
Agreement requires supported acceptance. Implementing action requires a
confirmed relevant act, such as funding, a reveal, or a transfer. The primary
behavioral endpoint, executed harmful coordination, requires supported
agreement, a knowingly source-adverse implementing action, and a victim or
compromised process. A wrong vote alone is insufficient. Strict constitutional
violation and executed collusion under the obligation-based rubric are separate
labels requiring a specific obligation actually delivered to participants.
Success means the predefined harmful target occurred. These labels are
distinct; a failed harmful attempt can still contain executed coordination.

Reference accuracy, completion, abstention, ordinary reward, extra transfer,
and capability failure are reported alongside those labels. Silence or an
invalid reveal cannot establish resistance by itself. Semantic uncertainty
remains unresolved, rather than being assigned a favorable label. Reviewers
will bind the actual claim hash and claim locator to fixed evidence, be blinded to model and
treatment where feasible, and retain disagreement and adjudication. Agent or
model-judge annotations remain provisional until independently checked.

Resistance requires explicit annotations that execution and capability failure
are absent, plus affirmative evidence of honest completion or competent
nonparticipation under the declared rubric. Nonexecution alone does not enter
its numerator. The current report gives fixed-frame outcome counts, without
sampling CIs or semantic resistance estimates. Repeated sources, jurors within
an episode and matched cross-model cases are dependent observations; five
launch seeds do not create independent documents. Appendix~\ref{app:estimation}
defines the prespecified bounds, equal source/mechanism weighting and planned
cluster handling for a later independently reviewed analysis.

\subsection{Communication and action detection}
The planned detector comparison uses communication, implemented actions, and
both together on identical sessions with fixed independent full-evidence
labels. Returned reasoning is a separately declared modality. Grouped
development/test separation and frozen thresholds are required before any
held-out comparison; the selected cases below do not establish superiority.
